%------------------------------------------------------------------------------%

\usepackage{calculo_varias_variaveis-1}

\title{Linearização}

%------------------------------------------------------------------------------%
\begin{document}

\addtitle

%------------------------------------------------------------------------------%
\section{Aproximação Linear}

%------------------------------------------------------------------------------%
\begin{frame}{Plano Tangente a superfície $z=f(x, y)$}
  Sabemos que o plano tangente à superfície \(z=f(x, y)\) \\
  de uma função diferenciável $f$
  no ponto \(\left(x_0, y_0, f(x_0, y_0)\right)\) é
  \[
    \emph{f_x(x_0, y_0)} (x-x_0) +
    \emph{f_y(x_0, y_0)} (y-y_0) -
    (z-z_0) = 0
  \]

  \pause
  Rearranjando
  \[
    z
    = \emph{f(x_0, y_0)}
    + \emph{f_x(x_0, y_0)} (x-x_0)
    + \emph{f_y(x_0, y_0)} (y-y_0)
  \]
\end{frame}

%------------------------------------------------------------------------------%
\begin{frame}{Aproximação Linear}
  Aproximação linear de uma função diferenciável de duas variáveis \\
  em torno do ponto \((x_0, y_0)\)
  \pause
  \[
    L(x, y)
    = \emph{f(x_0, y_0)}
    + \emph{f_x(x_0, y_0)} (x-x_0)
    + \emph{f_y(x_0, y_0)} (y-y_0)
  \]
\end{frame}

%------------------------------------------------------------------------------%
\addtocounter{exemplo}{1}
\begin{frame}{Exemplo \theexemplo}
  Encontre a linearização de
  \[
    f(x, y) = x^2 - xy + \frac{y^2}{2} + 3
  \]
  no ponto \((3, 2)\)
\end{frame}

\begin{frame}{Exemplo \theexemplo}
  \onslide<+->{Avaliando as derivadas parciais}
  \begin{align*}
    \onslide<+->{f_x(x, y)}
    \onslide<+->{& = \D{}{x} \left(x^2 - xy + \frac{y^2}{2} + 3\right)}
    \onslide<+->{= 2x - y \\[5mm]}
    \onslide<+->{f_y(x, y)}
    \onslide<+->{& = \D{}{y} \left(x^2 - xy + \frac{y^2}{2} + 3\right)}
    \onslide<+->{= -x + y}
  \end{align*}
\end{frame}

\begin{frame}{Exemplo \theexemplo}
\onslide<+->{Avaliando a função e as derivadas parciais no ponto \((3, 2)\)}
\begin{align*}
  \onslide<+->{f(3, 2)}
  \onslide<+->{& = \left(x^2 - xy + \frac{y^2}{2} + 3\right)\evalat{(3, 2)}{}}
  \onslide<+->{= 3^2 - 3\times 2 + \frac{2^2}{2} + 3}
  \onslide<+->{= 8 \\[3mm]}
  \onslide<+->{f_x(3, 2)}
  \onslide<+->{& = \left(2x - y\right)\evalat{(3, 2)}{}}
  \onslide<+->{= 2\times 3 - 2}
  \onslide<+->{= 4\\[3mm]}
  \onslide<+->{f_y(3, 2)}
  \onslide<+->{& = \left(-x + y\right)\evalat{(3, 2)}{}}
  \onslide<+->{= -3 + 2}
  \onslide<+->{= -1}
\end{align*}
\end{frame}

\begin{frame}{Exemplo \theexemplo}
  \onslide<+->{Aproximação linear}
  \begin{align*}
    \onslide<+->{L(x, y)}
    \onslide<+->{& = f(x_0, y_0) + f_x(x_0, y_0)(x-x_0) + f_y(x_0, y_0)(y-y_0) \\[3mm]}
    \onslide<+->{& = f(3, 2) + f_x(3, 2)(x-3) + f_y(3, 2)(y-2) \\[3mm]}
    \onslide<+->{& = 8 + 4(x-3) + (-1)(y-2) \\[3mm]}
    \onslide<+->{& = 4x -y -2}
  \end{align*}
\end{frame}

%------------------------------------------------------------------------------%
\section{Erro na Aproximação Linear}

%------------------------------------------------------------------------------%
\begin{frame}{Erro na Aproximação Linear}
  Se $f$ possui primeira e segunda derivadas parciais
  \pause
  sobre um conjunto aberto
  contendo um retângulo $R$ centrado em \((x_0, y_0)\)
  \pause
  e se $M$ é qualquer
  limitante superior para os valores de \(\abs{f_{xx}}\),  \(\abs{f_{yy}}\)
  e \(\abs{f_{xy}}\) em $R$
  \pause
  \[
    \emph{\abs{f_{xx}} \leq M} \qquad
    \emph{\abs{f_{yy}} \leq M} \qquad
    \emph{\abs{f_{xy}} \leq M} \qquad
    \forall (x, y) \in R
  \]
  \pause
  então o erro $E(x, y)$ na linearização satisfaz
  \[
    \emph{\abs{E(x, y)} \leq \frac{M}{2} \big(\abs{x-x_0}+\abs{y-y_0}\big)^2}
  \]
\end{frame}

%------------------------------------------------------------------------------%
\addtocounter{exemplo}{1}
\begin{frame}{Exemplo \theexemplo}
  Encontre um limitante superior para o erro na aproximação
  \[
    f(x, y) \approx L(x, y) = 4x -y -2
  \]
  no retângulo
  \[
    R: \abs{x-3} \leq 0,1 \,,\quad \abs{y-2} \leq 0,1
  \]
\end{frame}

\begin{frame}{Exemplo \theexemplo}
  Calculando os módulos das derivadas segundas
  \begin{align*}
    \onslide<+->{f_{xx}(x, y)}
    \onslide<+->{& = \DS{}{x} \left(x^2 - xy + \frac{y^2}{2} + 3\right) \\[3mm]}
    \onslide<+->{& = \D{}{x}\left(2x - y\right) \\[3mm]}
    \onslide<+->{& = 2 \\[3mm]}
    \onslide<+->{\abs{f_{xx}(x, y)}}
    \onslide<+->{& = 2}
  \end{align*}
\end{frame}

\begin{frame}{Exemplo \theexemplo}
  \begin{align*}
    \onslide<+->{f_{yy}(x, y)}
    \onslide<+->{& = \DS{}{y} \left(x^2 - xy + \frac{y^2}{2} + 3\right) \\[3mm]}
    \onslide<+->{& = \D{}{y}\left(-x + y\right) \\[3mm]}
    \onslide<+->{& = 1 \\[3mm]}
    \onslide<+->{\abs{f_{yy}(x, y)}}
    \onslide<+->{& = 1}
  \end{align*}
\end{frame}

\begin{frame}{Exemplo \theexemplo}
  \begin{align*}
    \onslide<+->{f_{xy}(x, y)}
    \onslide<+->{& = \DM{}{x}{y} \left(x^2 - xy + \frac{y^2}{2} + 3\right) \\[3mm]}
    \onslide<+->{& = \D{}{x}\left(-x + y\right) \\[3mm]}
    \onslide<+->{& = -1 \\[3mm]}
    \onslide<+->{\abs{f_{xy}(x, y)}}
    \onslide<+->{& = 1}
  \end{align*}
\end{frame}

\begin{frame}{Exemplo \theexemplo}
  \onslide<+->{Como o máximo das derivadas segundas em $R$ é 2}

  \onslide<+->{podemos escolher $M=2$, portanto}
  \begin{align*}
      \onslide<+->{\abs{E(x, y)}}
      \onslide<+->{& \leq \frac{M}{2} \big(\abs{x-x_0}+\abs{y-y_0}\big)^2 \\[3mm]}
      \onslide<+->{& = \frac{2}{2} \big(\abs{x-3}+\abs{y-2}\big)^2 \\[3mm]}
      \onslide<+->{& = \big(\abs{x-3}+\abs{y-2}\big)^2}
  \end{align*}
\end{frame}

\begin{frame}{Exemplo \theexemplo}
  \onslide<+->{Como
  \[
    \abs{x-3} \leq 0,1 \qquad \abs{y-2} \leq 0,1
  \]}
  \begin{align*}
    \onslide<+->{\abs{E(x, y)}}
    \onslide<+->{& \leq \big(\abs{x-3}+\abs{y-2}\big)^2\\[3mm]}
    \onslide<+->{& \leq \big(0,1+0,1\big)^2\\[3mm]}
    \onslide<+->{& = \big(0,2\big)^2\\[2mm]}
    \onslide<+->{& = 0,04}
  \end{align*}
\end{frame}

%------------------------------------------------------------------------------%
\section{Lista Mínima}

%------------------------------------------------------------------------------%
\begin{frame}{Lista Mínima}
  Cálculo Vol. 2 do Thomas 12\textsuperscript{a} ed. -- Seção 14.6
  \begin{enumerate}
    \item Estudar o texto da seção
    \item Resolver os exercícios: 20, 22, 24, 26, 30, 34, 36, 42, 48
  \end{enumerate}

  \vfill
  Atenção: A prova é baseada no livro, não nas apresentações
\end{frame}

%------------------------------------------------------------------------------%
\end{document}
%------------------------------------------------------------------------------%
